
\subsubsection{2.1 Semantic Entropy}

Farquhar et al. (2024) introduced semantic entropy as an uncertainty measure that accounts for semantic equivalence among generated responses. SE clusters responses using bidirectional natural language inference and computes entropy over cluster proportions. The method achieves AUROC 0.76-0.97 for hallucination detection across TruthfulQA, TriviaQA, and other benchmarks.

SE computation requires generating N=10-20 responses per query with temperature sampling, embedding each response, and performing O($N^{2})$ pairwise entailment checks using a DeBERTa-v3 NLI model \citep{He2021DeBERTa}. This computational overhead limits deployment in latency-sensitive applications.

\subsubsection{2.2 Semantic Entropy Probes}

Kossen et al. (2024) proposed Semantic Entropy Probes (SEPs) as lightweight classifiers trained on LLM hidden states to predict SE labels. SEPs extract hidden states from specific layers and token positions, then train logistic regression probes to map these representations to SE predictions. Published results report that SEPs achieve AUROC within 2-3\% of full SE on TruthfulQA while requiring only a single forward pass.

The SEP approach assumes that SE-relevant information is encoded in hidden states at particular layers and positions. The present work tested this assumption under one specific configuration.

\subsubsection{2.3 Related Uncertainty Estimation Approaches}

Kernel Language Entropy (KLE) constructs positive semidefinite kernels over response sets for uncertainty quantification \citep{Nikitin2024Kernel}. KLE demonstrates that similarity structure contains uncertainty-relevant information, which motivated the similarity augmentation approach tested in this work.

First-token entropy uses the entropy of the initial token distribution as an uncertainty proxy. Prior experiments in this research pipeline found rho = 0.13 correlation between first-token entropy and true SE, indicating limited effectiveness of token-level metrics for capturing semantic uncertainty.
